\section{Shared Resources}

\subsection{Race Conditions}

\begin{definition}

                A \textbf{Race condition} occurs in concurent programming when the correctness of the system depends on
                the specific interleaving or ordering of execution of operations.
            
\end{definition}

\begin{definition}

                A \textbf{Low-level race condition} is a race condition that occurs because of insufficient
                synchronization of memory access (Data Races). While the \textbf{High-level race condition} is a race
                condition that
                occurs because of a wrong decomposition/algorithm. The program might ensure mutual exclusion on the
                shared resources but bad interleavings still lead to wrong results.
            
\end{definition}

            Generally we want one of three things when it comes to shared ressources:
            \begin{itemize}
    \item \textbf{Thread local:} We dont use the ressource of other threads.
    \item \textbf{Read only (Immutable):} We only read and do not change the ressource
    \item \textbf{Mutual exclusion (Synchronization):} We use the ressource of other threads but we ensure that
                    only one thread accesses it at a time.
\end{itemize}

